\section{Commitments and the single-exponential bound}\label{sec:commitments}

\subsection{The relative-boundary commutation lemma}

\begin{lemma}[Commuting disjoint supports]\label{lem:commute}
Let $a,b$ be simultaneously legal formal replacements in a triangulation,
and suppose that their sets of consumed tetrahedron incarnations are disjoint.
After applying $a$, the transported event $b$ remains legal; the symmetric
statement holds. The results of $ab$ and $ba$ have a canonical combinatorial
isomorphism fixing every surviving old tetrahedron and respecting the two
replacement regions, their boundary attachments and their transported cocycles.
\end{lemma}
\begin{proof}
Regard each consumed region as a separate formal bipyramid before the global
face identifications. Both fillings have the same labelled six-facet boundary.
One can replace both fillings simultaneously and then impose the original
attaching maps. Replacing them in either order gives exactly this quotient.
This also covers the case that a boundary face of one region is attached to
a boundary face of the other: the common attaching permutation is retained
in the simultaneous construction.

The internal face of a $2$--$3$ event belongs to its two consumed tetrahedra.
The complete central-edge star of a $3$--$2$ event belongs to its three
consumed tetrahedra. A disjoint replacement changes neither of these internal
configurations. Its updates to exterior attachments merely apply the same
relative-boundary identifications. Hence the second event stays legal.

For the marking, the five heights in each formal region are reconstructed
from the same old edge differences. Changes of additive row constants have
no effect. The regions use disjoint old tetrahedron interiors and retain the
same boundary edge values. The two independently obtained sets of new rows
therefore agree after the stated isomorphism and normalization. No claim
that the intermediate coherent surfaces have identical topology is required.
\end{proof}

Disjoint support is sufficient, not necessary. Overlapping moves can also
have special relations, as in the earlier Pachner-graph literature
\cite{BurtonPachner}; the present proof deliberately requires only the
uniformly checked relation in Lemma~\ref{lem:commute}.

\subsection{A constant alphabet on each incarnation}

Assign an incarnation one of the following eleven symbols:
\[
 \mathcal A=\{\text{idle}\}
 \cup\{\text{down at edge }e:0\le e<6\}
 \cup\{\text{up at face }f:0\le f<4\}.
\]
The six edges use a fixed ordering of the pairs of local vertices.
An enabled downward event is \emph{ready} if all three source incarnations
are committed to their local copy of its central edge. An upward event is
ready if both source incarnations are committed to the common face. Upward
events are disabled when the total upward budget is exhausted.

\begin{lemma}[Ready events cannot compete]\label{lem:ready}
Two different ready events consume disjoint sets of incarnations.
\end{lemma}
\begin{proof}
A shared incarnation would have to select the same symbol for both events.
The idle symbol selects no event. A selected local face identifies its unique
partner and hence its unique $2$--$3$ event. A selected local edge identifies
its entire global edge star and hence the unique legal $3$--$2$ event at
that edge. Thus the purported events would be equal.
\end{proof}

The committed search uses a fixed deterministic order of ready events. If
one exists, it executes the first, preserves commitments on survivors and
leaves new incarnations unassigned. If none exists, it assigns the first
unassigned live incarnation one of the alphabet symbols, branching on every
choice. If none is unassigned, that branch stops. An assignment need not be
consistent with any eventual move; inert assignments are allowed and counted.

The implementation evaluates a state predicate at each newly encountered
marked endpoint, before deciding whether to continue. This avoids the extra
assumption that a useful surface survives every later downward choice.

\begin{theorem}[Complete commitment encoding]\label{thm:commitment}
Let $T$ have $t$ tetrahedra and let $U\ge0$. Every permitted trace defining
an endpoint in $\Reach_U(T,h)$ is equivalent to an execution of the
deterministic committed scheduler.
At most $3t+5U$ incarnation symbols are needed. Consequently the number of
commuting trace classes, including classes of all finite prefixes, is at most
\begin{equation}\label{eq:commitment-count}
 11^{3t+5U}.
\end{equation}
For $U=0$, the alphabet can be restricted to seven symbols and the corresponding
bound is $7^{3t}$.
\end{theorem}
\begin{proof}
Fix a target finite trace. For each incarnation appearing in it, record the
local port of the event that eventually consumes it; if it survives to the
target endpoint, record idle. Follow these assignments in the deterministic
scheduler, interpreting them under the transported local labels.

Suppose an event is currently ready. Consider the first event in the
remaining target trace that touches one of its source incarnations. Until
that first touch, the ready event's common face or full central-edge star is
unchanged. The touched incarnation's commitment specifies the local port of
its eventual consuming event. This is the ready event's port, so uniqueness
at that port implies that the first touching event is exactly the ready
event. All earlier target events have disjoint consumed supports. By
Lemma~\ref{lem:commute}, the ready event can be moved to the front across all
of them. This interchange preserves the endpoint, the marking and the total
upward count. Its new prefix cannot exceed the total allowed upward count.

If no event is ready, assign the next unassigned incarnation its recorded
symbol. Once the participants of the first remaining target event have
received their symbols, that event is ready. Thus this process cannot stop
before all target events have been performed, apart from harmless
reordering. All target survivors are idle, so it stops at the required
endpoint after the assignments are completed. This proves completeness and
trace equivalence without asserting confluence of uncommitted downward
rewriting.

If the execution has $u$ upward and $d$ downward events, the total number of
incarnations is
\[
 N=t+3u+2d\le t+3u+2(t+u)=3t+5u\le3t+5U.
\]
Each is assigned at most once. Encode the assignments in the deterministic
order in which the scheduler requests them, and pad a shorter encoding to
length $3t+5U$ by idle symbols. Replaying a word determines one execution.
If two trace classes had the same encoding, they would both commute to
that same execution and would be equal. Hence there are at most the number
of such words. The same construction applied to any target prefix, giving
its survivors idle, proves the assertion for all prefix classes. Omitting
upward symbols gives the last statement.
\end{proof}

\begin{remark}[Assignments are not geometric steps]
The encoding algorithm is a proof device as well as a reference implementation.
Many assignments lead nowhere. Its tree has at most
$(11^{N+1}-1)/10$ assignment vertices, with polynomially many forced-move
operations along an assignment path. It can be much slower on a small
triangulation than direct trace reduction. The bound does not justify a
claim of practical speed by itself.
\end{remark}

\subsection{Sleep sets give an efficient realization of the same bound}

At a search node maintain a set of sleeping event identities. Explore the
enabled events in deterministic order. Skip a sleeping event. After exploring
an event $a$, add $a$ to the local sleep set for later sibling branches.
For a child reached by $b$, retain only sleeping events whose consumed support
is disjoint from $b$'s support. Thus an event is remembered across a replacement
exactly when the proved commutation rule permits it to move past that replacement.

\begin{theorem}[Sleep-set coverage and uniqueness]\label{thm:sleep}
With exact incarnation identities, one canonical event per eligible edge or
face, and the independence relation of Lemma~\ref{lem:commute}, sleep-set
exploration visits at most one representative of each commuting trace class.
Every reachable marked endpoint has a visited representative. Hence the
number of visited trace nodes is at most $11^{3t+5U}$, or $7^{3t}$ when $U=0$.
\end{theorem}
\begin{proof}
For uniqueness, suppose equivalent traces first diverge at a state where the
earlier branch takes $a$ and a later branch takes $b$. The later branch starts
with $a$ sleeping. Legal interchanges preserve exact event identities, so it
must eventually perform that same event $a$. None of its consumed
incarnations can have been destroyed in the intervening prefix: an
incarnation never reappears after destruction. Every intervening event thus
has consumed support disjoint from $a$. Initially $a$ is enabled, and
Lemma~\ref{lem:commute} inductively keeps it enabled across that prefix.
The sleep-update rule retains $a$ throughout, so its eventual occurrence is
forbidden. This is a contradiction.

For coverage, choose among the finitely many representatives of any target
trace class the first one in the deterministic search order. If it were
blocked by a sleeping event, the earlier occurrence that put that event to
sleep could be commuted across the intervening independent events, giving
an earlier representative of the same class. This contradicts the choice.
Equivalently, repeatedly move the earliest blocking event left until no
blocking event remains; the finite lexicographic order decreases. Thus
the chosen representative is visited. Lemma~\ref{lem:commute} supplies
the persistence of all transported enabled events used in this argument.

The total-upward restriction is invariant under these interchanges. Protected
initial cells, when present, are also invariant. Apply
Theorem~\ref{thm:commitment} to the resulting representatives.
\end{proof}

This proof uses the standard sleep-set principle \cite{Godefroid}, but its
hypotheses are checked for these geometric events. The implementation does
not merge search continuations merely because their raw triangulations agree.
Different remaining upward budgets and sleep sets can permit different
continuations. A separate exact state table suppresses repeated endpoint
queries only; exploration continues from every scheduled state.

Independence here concerns simultaneously enabled events, or legal adjacent
interchanges. It is not a static relation inferred only from the consumed
sets of arbitrary historical events: a later event can consume an output
born from an earlier event even though their consumed sets are disjoint.
The persistence and nonresurrection argument above is what justifies the
actual sleep update.

\subsection{A genuine factorial-to-subset separation}

Suppose $p$ independent downward events are enabled, no other downward event
is created, and each can be executed at most once. Unreduced all-prefix
enumeration visits
\begin{equation}\label{eq:factorial}
 N_{\rm naive}(p)=\sum_{j=0}^{p}\frac{p!}{(p-j)!}
 =p!\sum_{k=0}^{p}\frac1{k!}.
\end{equation}
Every subset of events is a distinct prefix endpoint and all orderings of
that subset commute. Sleep exploration visits exactly $2^p$ nodes. Full
maximal orders alone are reduced from $p!$ to one.

The artifact realizes this situation by disjoint inverse bipyramids in a
layered solid torus and verifies the move inventory and endpoints. It is
therefore an actual geometric control. Equation~\eqref{eq:factorial} is a
proof under its stated event-inventory hypothesis; finite tests of the
construction are reported at the tested sizes. We do not extrapolate the
measured construction to a universal claim about all solid-torus
triangulations.
